\begin{abstract}

Pairwise sequence identity from haplotype-resolved assemblies directly estimates $1 - 2\mu t$, the theoretically optimal quantity for genetic relatedness. We present HPRCv2-IBD, a Rust toolkit that performs IBD detection, local ancestry inference, and relatedness estimation directly from pangenome alignments by applying Hidden Markov Models to sequence identity computed via \texttt{impg}. Validation on 223 Human Pangenome Reference Consortium individuals across chromosomes 10, 11, 12, and 20 uses matched reference panels and gold-standard VCF-based tools (hap-ibd, RFMix). For IBD, 11/11 (100\%) hap-ibd pairs rank in the top 10\% by robust statistics across six genomic regions, with 5/6 regions achieving perfect precision--recall at calibrated thresholds. The identity gap between IBD and non-IBD pairs ($\Delta\mu \approx 0.003$) is reproducible across all chromosomes. Eight independent algorithmic improvements converge to the same performance ceiling, establishing a measurement rather than modeling limitation. For local ancestry, a 3-way (EUR/AFR/AMR) pairwise contrast model achieves 76.45\% per-haplotype concordance with RFMix across 100 comparisons; the dominant ancestry reaches F1~=~0.815 while minority ancestries reflect the fundamental per-window signal-to-noise limitation. Simulation with known ground truth achieves 97.95\% concordance, exceeding RFMix (95.9\%) by 2.05 percentage points; simulated IBD reaches pair-level F1~=~0.514, competitive with hap-ibd (0.489). Grandparent painting in the Platinum pedigree (CEPH 1463) achieves 99.53\% mean accuracy from assembly alignments alone. These results establish pangenome identity as the most direct route to the TMRCA signal that coalescent theory identifies as optimal.

\end{abstract}
